% \iffalse meta-comment
%
% Copyright (C) 2017- by Yanshuo Chu <yanshuoc@gmail.com>
%
% This file may be distributed and/or modified under the
% conditions of the LaTeX Project Public License, either version 1.3a
% of this license or (at your option) any later version.
% The latest version of this license is in:
%
% http://www.latex-project.org/lppl.txt
%
% and version 1.3a or later is part of all distributions of LaTeX
% version 2004/10/01 or later.
%
% \fi
%
% \section{报告首页：深圳本科}
%
% ^^A ======================================================================
% ^^A Module report-cover-shenzhen-bachelor: 深圳本科开题与中期的首页表单
% ^^A ======================================================================
%
% \changes{v3.2a}{2026/09/13}{深圳本科那两份报告的首页，中英各一套}
% 基准那一份在 \file{src/pages/hit-report-cover.dtx}。这一份是另一页，与基准
% 没有重叠的部件，所以整页自己排，登记成 \texttt{report-cover-page} 的
% 深圳本科那一档。
%
% \emph{数取自 iota-hit 的 \texttt{src/pages/report/shenzhen/cover.typ}。} 那边是
% 拆开 2026 年 9 月版两份原件（深圳本科毕业论文（设计）开题报告、同版中期报告）
% 的 \texttt{.docx} 量出来的。本仓库没有原件，照抄，没有复核。
%
% \emph{整页照中期那一份排，只换字。} 开题那一份原件的线是一串带 \texttt{w:u}
% 的手打半角，表还是浮动的，一行一个长短；中期那份用的是真表格加单元格边框，
% 齐整得多。两份的列宽、空段、表在页上的落点共用一套，开题与中期的差别只剩
% 校名下那一行与题目那一格的标签，跟词条走。
%
% \emph{与本部本科那一页的区别：}
% \begin{itemize}
% \item 头两行是\emph{文字}，小初华文新魏加粗；本部那页是校名题字的矢量图。
% \item 字段是一张两列的表，左列标签、右列的下划线是单元格的下边框；本部那页
%   是五段，一段一格、首行缩进。
% \item 正文照论文排，见 \file{src/config/hit-layout.dtx}。
% \end{itemize}
%
%    \begin{macrocode}
%<*report-cover-shenzhen-bachelor>
\bool_lazy_all:nT
  {
    { ! \bool_if_p:N \g__hit_final_bool }
    { \bool_if_p:N \g__hit_shenzhen_bool }
    { \bool_if_p:N \g__hit_bachelor_bool }
  }
  {
%    \end{macrocode}
%
% \emph{整页自成一节}，上下 72\,bp、左右 90\,bp，不出页眉页脚，见
% \file{src/config/hit-layout.dtx} 的 \option{cover} 那一档部件差额。
% 表在页上\emph{整体居中}（开题原件是浮动表 \texttt{tblpXSpec="center"}，
% 中期是 \texttt{w:jc="center"}），所以只认列宽，不认页边距。
%
% \emph{列宽照原件的 \texttt{w:tblGrid}。} 中文版两个阶段各一套（开题
% 2250／3691 缇，中期 2206／3735，差 2.2\,bp，同一次改版发的两份表多半是手拉的，
% 跟原件不归一）；英文版两个阶段同一套 3357／4961，比中文那两份宽，标签那一列
% 要装下双语。单元格左右留白 108 缇，是 Word 没写 \texttt{w:tblCellMar} 时的默认值。
%    \begin{macrocode}
    \dim_new:N \l__hit_szb_label_dim
    \dim_new:N \l__hit_szb_value_dim
    \dim_new:N \l__hit_szb_inset_dim
    \dim_new:N \l__hit_szb_indent_dim
    \dim_new:N \l__hit_szb_topskip_dim
    \dim_set:Nn \l__hit_szb_inset_dim { 5.4bp }
    \bool_if:NTF \g__hit_en_bool
      {
        \dim_set:Nn \l__hit_szb_label_dim { 167.85bp }
        \dim_set:Nn \l__hit_szb_value_dim { 248.05bp }
      }
      {
        \bool_if:NTF \g__hit_proposal_bool
          {
            \dim_set:Nn \l__hit_szb_label_dim { 112.5bp }
            \dim_set:Nn \l__hit_szb_value_dim { 184.55bp }
          }
          {
            \dim_set:Nn \l__hit_szb_label_dim { 110.3bp }
            \dim_set:Nn \l__hit_szb_value_dim { 186.75bp }
          }
      }
%    \end{macrocode}
%
% \emph{标签摊开。} 摊的只是核心汉字之间的缝，核心指头一个汉字到末一个汉字
% 那一截，后面的冒号不摊。
%
% 中文版钉的是\emph{缝}不是总宽：缝共 $(5.5-\text{核心汉字数})$ 个小三格，
% 匀给核心里的每一道缝；汉字跟本行的字号走。所以题目那一行（小二）比另六行
% （小三）宽 9\,bp，原件量的正是 106.5 与 97.5。
%
% 英文版的缝是「几个半角」，按本行汉字字号折算，题目 3 个、六格 4 个，而且
% 只有两个字的标签才摊：原件里「指导教师」四个字是挨着的，一个空格都没敲。
%    \begin{macrocode}
    \int_new:N \l__hit_szb_idx_int
    \tl_new:N \l__hit_szb_head_tl
    \tl_new:N \l__hit_szb_tail_tl
    \seq_new:N \l__hit_szb_parts_seq
    \cs_new_protected:Npn \hit@szb@spread:nnn #1#2#3
      {
        \int_zero:N \l__hit_szb_idx_int
        \tl_map_inline:nn {#3}
          {
            \int_incr:N \l__hit_szb_idx_int
            \int_compare:nT { 1 < \l__hit_szb_idx_int <= #2 }
              { \skip_horizontal:n {#1} }
            ##1
          }
      }
    \cs_generate_variant:Nn \hit@szb@spread:nnn { nnV }
%    \end{macrocode}
%
% \emph{英文版标签的尾巴。} 这一族词条都写成「中文/ English：」，拿斜杠切开：
% 前一半是核心汉字，跟本行字号走；后一半连着斜杠一起排 Times 加粗三号，
% 末尾那个全角冒号是宋体、不加粗，原件量出来正是这样。
%    \begin{macrocode}
    \cs_new_protected:Npn \hit@szb@latin:n #1
      {
        \seq_set_split:Nnn \l__hit_szb_parts_seq { ： } {#1}
        \fontsize { 16bp } { 16bp } \selectfont
        \int_compare:nTF { \seq_count:N \l__hit_szb_parts_seq > 1 }
          {
            { \rmfamily \bfseries / \seq_item:Nn \l__hit_szb_parts_seq { 1 } }
            { \songti ： }
          }
          { \rmfamily \bfseries / #1 }
      }
    \cs_generate_variant:Nn \hit@szb@latin:n { V }
%    \end{macrocode}
%
% \meta{词条名}\meta{核心汉字数}\meta{英文版的缝是几个半角}\meta{英文版汉字字号}。
% 后两个参数中文版用不上。词条名与内封、与本部那两页同名，一个新键都不加。
%    \begin{macrocode}
    \tl_new:N \l__hit_szb_term_tl
    \cs_new_protected:Npn \hit@szb@label:nnnn #1#2#3#4
      {
        \tl_set:Nx \l__hit_szb_term_tl
          {
            \bool_if:NTF \g__hit_en_bool
              { \use:c { hit@term@cover@#1@label@shenzhen@bachelor@en } }
              { \use:c { hit@term@cover@#1@label@shenzhen@bachelor@zh } }
          }
        \bool_if:NTF \g__hit_en_bool
          {
            \seq_set_split:NnV \l__hit_szb_parts_seq { / } \l__hit_szb_term_tl
            \tl_set:Nx \l__hit_szb_head_tl
              { \seq_item:Nn \l__hit_szb_parts_seq { 1 } }
            \tl_set:Nx \l__hit_szb_tail_tl
              { \seq_item:Nn \l__hit_szb_parts_seq { 2 } }
            {
              \fontsize {#4} {#4} \selectfont
              \int_compare:nTF { #2 = 2 }
                {
                  \hit@szb@spread:nnV { \fp_to_dim:n { #3 * #4 / 2 } } {#2}
                    \l__hit_szb_head_tl
                }
                { \l__hit_szb_head_tl }
            }
            \tl_if_empty:NF \l__hit_szb_tail_tl
              { { \hit@szb@latin:V \l__hit_szb_tail_tl } }
          }
          {
            \hit@szb@spread:nnV
              { \fp_to_dim:n { ( 5.5 - #2 ) * 15bp / ( #2 - 1 ) } } {#2}
              \l__hit_szb_term_tl
          }
      }
%    \end{macrocode}
%
% \emph{一行的形状。} 标签那一格右留白不占位，让装不下的标签探出去而不是折行：
% 原件里 Word 也是这么办的，英文版「指导教师/ Supervisor：」量得 159.2，
% 比格子内缘 157.05 还宽，探出去没折。
%
% \emph{值在格里居中。} 原件里那一格的字是居中的，不是顶着左缘排；居中算的是
% \emph{整格宽}，不减单元格的左右留白（那个留白只管标签那一格）。逐字量 iota
% 的渲染件：值那一格宽 184.55、左缘 261.61，「张三」两个字三十宽，居中落在
% 338.89，一个数不差。
%
% \emph{格里每一行画一条线，格自己长高。} 值写长了就在格里折行，折出来的每一行
% 底下都有一条线，下面几格跟着往下让。原件那一格里本来就是一串带 \texttt{w:u}
% 的空格，填了字就是字的下划线，所以线跟着行走，不是单元格边框。
%
% \emph{折出来的行按「行盒加段前」排，不按行距排。} 这一处\emph{不照原件}：
% Word 里段内换行没有段前，折出来的两条线只隔一个行盒，比格与格之间的步进挤
% 一截，填表的人看着像两种线距。加上段前之后整页的线等距，折了行的那一格看着
% 就是多了一格。iota 也是这么定的。
%
% \emph{行盒按 Word 算。} 一格一段，段前 12\,bp、行距 1.5 倍，行盒高取标签与值
% 里高的那个；英文版标签带一截 16\,bp 的尾巴，六格的步进因此是 43.13 而不是
% 41.18。这一档由 \option{report-cover-field} 给。
%
% \emph{整格一种字，线跟着这一格挑。} 有汉字的格走中文那一档：Word 画的线心在
% 基线下「一个 em 除以 8 加 0.5\,bp」、线宽 $12/256$ 个 em，拿中易三副字各排
% 十六个字号量出来的，不照字体的 \texttt{post} 表。一个汉字都没有的格（学号那种
% 纯数字、空着的格）走 Times 那一档，照它 \texttt{post} 表的 $-223/2048$ 与
% $100/2048$；基线也跟着高一截，行盒矮一截。
%
% \emph{与 iota 的差别记在这里：} 那边是\emph{逐行}判有没有汉字，中西混排折了行
% 时每一行各挑各的。这里按整格判，单行的格（六格填的都是这一路）完全一样，
% 中西混着又折了行的格会整格取一档。
%    \begin{macrocode}
    \tl_new:N \l__hit_szb_font_tl
    \str_new:N \l__hit_szb_scan_str
    \bool_new:N \l__hit_szb_cjk_bool
    \box_new:N \l__hit_szb_cell_box
    \dim_new:N \l__hit_szb_pitch_dim
    \dim_new:N \l__hit_szb_ascent_dim
    \dim_new:N \l__hit_szb_size_dim
    \dim_new:N \l__hit_szb_label_ascent_dim
    \dim_new:N \l__hit_szb_center_dim
    \dim_new:N \l__hit_szb_thick_dim
%    \end{macrocode}
%
% \emph{这一格是不是汉字那一档，看它填了什么。} 把值展开成字符串逐字看码位，
% $\ge$ U+2E80 的算汉字那一路。学号填数字就走西文、填汉字就走中文，不必另标。
%
% \emph{日期那一格例外，按语种定。} \cs{hit@info@date@zh} 一族是
% \texttt{protected} 的，\cs{edef} 展不开，\cs{text\_purify:n} 又把不认识的直接
% 丢掉；而「年月日」正是它排的时候才加的，从取值上根本看不见。中文那一档恒有
% 这三个字，英文那一档恒没有，所以直接按语种给。
%
% \emph{英文版没填英文就退回中文那一份}，与 iota 的 \texttt{value-of} 一致；
% 退回去之后这一格自然又是汉字那一档。
%
% \emph{填没填要看得见的那个存储。} 题目那一路是个派发宏
% （\cs{hit@info@title@cover@en} 恒展开成 \cs{hit@title@cover}\marg{en}），
% 空不空从它身上看不出来，所以另给一个探测用的名字：题目探
% \cs{hit@info@title@first@en}，六格探它自己那一份。
%    \begin{macrocode}
    \tl_new:N \l__hit_szb_value_tl
    \cs_new_protected:Npn \hit@szb@cjk@scan:n #1
      {
        \bool_set_false:N \l__hit_szb_cjk_bool
        \str_set:Nx \l__hit_szb_scan_str {#1}
        \str_map_inline:Nn \l__hit_szb_scan_str
          {
            \int_compare:nT { `##1 > "2E7F }
              { \bool_set_true:N \l__hit_szb_cjk_bool \str_map_break: }
          }
      }
    \cs_new_protected:Npn \hit@szb@pick:nnn #1#2#3
      {
        \tl_set:Nn \l__hit_szb_value_tl { \use:c {#3} }
        \bool_if:NT \g__hit_en_bool
          {
            \tl_if_blank:nF {#1}
              {
                \tl_if_empty:cF {#1}
                  { \tl_set:Nn \l__hit_szb_value_tl { \use:c {#2} } }
              }
          }
        \hit@szb@cjk@scan:V \l__hit_szb_value_tl
      }
    \cs_new_protected:Npn \hit@szb@pick:nn #1#2
      { \hit@szb@pick:nnn {#1} {#1} {#2} }
    \cs_generate_variant:Nn \hit@szb@cjk@scan:n { V }
%    \end{macrocode}
%
% 线心与线宽，以及基线在行盒顶下方多少。三个都按这一格挑的那一档算。
% 字号先存进寄存器：\cs{dimexpr} 里小数系数后面只能跟寄存器，
% 写成 \texttt{1.008 18bp} 那样的字面量会静默算错。
%    \begin{macrocode}
    \cs_new_protected:Npn \hit@szb@metrics:n #1
      {
        \dim_set:Nn \l__hit_szb_size_dim {#1}
        \bool_if:NTF \l__hit_szb_cjk_bool
          {
            \dim_set:Nn \l__hit_szb_center_dim
              { \l__hit_szb_size_dim / 8 + 0.5bp }
            \dim_set:Nn \l__hit_szb_thick_dim
              { \l__hit_szb_size_dim * 12 / 256 }
            \dim_set:Nn \l__hit_szb_ascent_dim
              {
                \l_docxlike_format_ascent_tl \l__hit_szb_size_dim
              }
          }
          {
            \dim_set:Nn \l__hit_szb_center_dim
              { \l__hit_szb_size_dim * 273 / 2048 }
            \dim_set:Nn \l__hit_szb_thick_dim
              { \l__hit_szb_size_dim * 100 / 2048 }
            \dim_set:Nn \l__hit_szb_ascent_dim
              {
                \c_docxlike_format_ascent_latin_tl \l__hit_szb_size_dim
              }
          }
      }
%    \end{macrocode}
%
% 一格若干行的线。画在零宽盒里，不占位、不参与折行；每一条比上一条低一个步进。
% 抬的量先算进一个寄存器再取值：\cs{raisebox} 与 \cs{rule} 一样收不下
% \cs{dimexpr}（它要给参数补单位），传进去静默变成零；而且它只把负号加在头一项
% 上，\texttt{-a+b} 不是 \texttt{-(a+b)}。
% 线画成上下各半个线宽的 \cs{vrule}，再整个抬到线心那一处，这样抬的量就是线心
% 本身。\cs{rule} 收不下 \cs{dimexpr}（它要给参数补单位），线宽会静默变成零。
%    \begin{macrocode}
    \dim_new:N \l__hit_szb_drop_dim
    \cs_new_protected:Npn \hit@szb@drop:nn #1#2
      {
        \dim_set:Nn \l__hit_szb_drop_dim { 0pt - ( #1 ) }
        \raisebox { \dim_use:N \l__hit_szb_drop_dim }
          [ \c_zero_dim ] [ \c_zero_dim ] {#2}
      }
    \cs_new_protected:Npn \hit@szb@rules:n #1
      {
        \hbox_overlap_right:n
          {
            \int_step_inline:nn {#1}
              {
                \hbox_overlap_right:n
                  {
                    \hit@szb@drop:nn
                      {
                        \l__hit_szb_ascent_dim + \l__hit_szb_center_dim
                        + \int_eval:n { ##1 - 1 } \l__hit_szb_pitch_dim
                      }
                      {
                        \vrule ~ width ~ \l__hit_szb_value_dim
                          ~ height ~ \dim_eval:n { 0.5 \l__hit_szb_thick_dim }
                          ~ depth ~ \dim_eval:n { 0.5 \l__hit_szb_thick_dim }
                        \scan_stop:
                      }
                  }
              }
          }
      }
%    \end{macrocode}
%
% \emph{把值排进一个 \cs{vtop}，行距锁成步进。} 字体要在锁之前发完：
% \cs{xiaoer} 这类字号命令自己会重设 \cs{baselineskip}，写在后面就把锁顶掉了，
% 实测题目折出来的第二行跑到 23.4\,bp 处（\cs{xiaoer} 的行距），不是 47.02。 \cs{vtop} 的基线是首行的，
% 与左边标签那一格对得齐；\cs{lineskiplimit} 调到负无穷，行高再大也不会让
% \hologo{TeX} 改用 \cs{lineskip}，行与行就恒是一个步进。行数由 \cs{prevgraf}
% 数，它是刚排完那一段的行数，准。
%    \begin{macrocode}
    \int_new:N \l__hit_szb_lines_int
    \cs_new_protected:Npn \hit@szb@fill:nnn #1#2#3
      {
        \vbox_set_top:Nn \l__hit_szb_cell_box
          {
            \dim_set_eq:NN \hsize \l__hit_szb_value_dim
            \dim_zero:N \parindent
            \dim_zero:N \parskip
            \int_set:Nn \tolerance { 9999 }
            \centering
            #2
            \dim_set_eq:NN \baselineskip \l__hit_szb_pitch_dim
            \dim_set:Nn \lineskiplimit { -\maxdimen }
            \noindent #3 \par
            \int_gset:Nn \g__hit_szb_prevgraf_int { \tex_prevgraf:D }
          }
        \int_set:Nn \l__hit_szb_lines_int
          { \int_max:nn { \g__hit_szb_prevgraf_int } {#1} }
      }
    \int_new:N \g__hit_szb_prevgraf_int
%    \end{macrocode}
%
% \emph{一整行。} 段前发在行外，行本身是个定高的 \cs{vbox}：高 $= (行数-1)$ 个
% 步进加一个行盒，正好让下一行的段前接上去之后，行与行的步进是整数个步进。
%
% \emph{标签与值各自从格顶往下摆，不共用一条基线。} 两边的盒子都压成零高零深，
% 再各按自己的上升部沉下去。标签走中文那一档、按整行的行盒算；值走这一格挑的
% 那一档。一格填纯西文时基线因此比标签高一截——原件填字后量的正是这样：题目
% 填英文的基线比「题目：」高 1.4，填中文与它齐。
%
% \meta{目标}\meta{值的字号}\meta{最少几行}\meta{定汉字标记的那一句}\meta{标签}
% \meta{值的字体}\meta{值}。
%    \begin{macrocode}
    \cs_new_protected:Npn \hit@szb@line:nnnnnnn #1#2#3#4#5#6#7
      {
        \par
        \docxlike_format_compute:n {#1}
        \dim_set:Nn \l__hit_szb_pitch_dim
          { \l_docxlike_format_lead_dim + \l_docxlike_format_before_dim }
        #4
        \hit@szb@metrics:n {#2}
        \dim_set:Nn \l__hit_szb_label_ascent_dim
          {
            \l_docxlike_format_ascent_tl \l_docxlike_format_size_dim
          }
        \docxlike_format_apply_font:nN {#1} \l__hit_szb_font_tl
        \hit@szb@fill:nnn {#3} { \l__hit_szb_font_tl #6 } {#7}
        \skip_vertical:n { \l_docxlike_format_before_dim }
        \nointerlineskip
        \hbox_to_wd:nn { \linewidth }
          {
            \skip_horizontal:n
              { \dim_eval:n { \l__hit_szb_indent_dim + \l__hit_szb_inset_dim } }
            \vbox_to_ht:nn
              {
                \dim_eval:n
                  {
                    \int_eval:n { \l__hit_szb_lines_int - 1 } \l__hit_szb_pitch_dim
                    + \l_docxlike_format_lead_dim
                  }
              }
              {
                \hbox:n
                  {
                    \hit@szb@drop:nn { \l__hit_szb_label_ascent_dim }
                      {
                        \hbox_to_wd:nn
                          { \dim_eval:n { \l__hit_szb_label_dim - \l__hit_szb_inset_dim } }
                          { #5 \hfil }
                      }
                    \hit@szb@rules:n { \l__hit_szb_lines_int }
                    \hit@szb@drop:nn { \l__hit_szb_ascent_dim }
                      { \hbox:n { \box_use:N \l__hit_szb_cell_box } }
                  }
                \vfil
              }
            \hfil
          }
        \par
      }
    \cs_new_protected:Npn \hit@szb@row:nnn #1#2#3
      {
        \hit@szb@line:nnnnnnn { report-cover-field } { 15bp } { 1 }
          { \hit@szb@pick:nn {#3} {#2} }
          {
            \heiti \xiaosan
            \hit@szb@label:nnnn {#1} { \hit@szb@core:n {#1} } { 4 } { 15bp }
          }
          { \songti \xiaosan }
          { \l__hit_szb_value_tl }
      }
    \cs_new:Npn \hit@szb@core:n #1
      { \str_if_eq:nnTF {#1} { supervisor } { 4 } { 2 } }
    \cs_new_protected:Npn \hit@szb@date@row:
      {
        \hit@szb@line:nnnnnnn { report-cover-field } { 15bp } { 1 }
          { \bool_set:Nn \l__hit_szb_cjk_bool { ! \bool_if_p:N \g__hit_en_bool } }
          {
            \heiti \xiaosan
            \hit@szb@label:nnnn { date } { 2 } { 4 } { 15bp }
          }
          { \songti \xiaosan }
          {
            \bool_if:NTF \g__hit_en_bool
              { \hit@info@date@en } { \hit@info@date@zh }
          }
      }
%    \end{macrocode}
%
% \emph{勾选框与本科内封共用一副}，见 \file{src/pages/hit-checkbox.dtx}。
%
% \emph{整页。} 头一行是 2026 年 9 月版新加的勾选行，按 \option{form} 勾一个：
% 小四、段前 12\,bp、行盒 19\,bp 固定值，两组之间空两个半角。原件那一段的
% \texttt{w:rPr} 写着小一，可那是\emph{段落标记}的字号，两个框与四个字的 run
% 一个 \texttt{sz} 都没写，跟 Normal 走。
%
% 接着一个空段：中文版二号、行距 1.25 倍，英文版小一、行盒 19\,bp 固定值，
% 两版都带 12\,bp 段前。这里按「段前加行高」折成一段竖距，不真排一个空段。
%
% 再是校名与文种。中文版两行都是小初华文新魏加粗、字距 2\,bp、字宽压到 90\%、
% 行距 1.2 倍。
%
% 英文版的校名与文种各排两行，上面那一行印中文、下面那一行印英文，
% 两行之间还有一个 8\,bp 的空段。英文版的文种中文那一行改成宋体小一加粗、
% 字宽不压、字距 4\,bp（它的 run 写着 \texttt{w:spacing="40"}）；英文那两行是
% Times 加粗，校名 20\,bp、文种 16\,bp，都不带字距——那个 20 写在\emph{段落标记}
% 上，run 上一个都没有。
%
% \emph{字宽压的是字身，字距不跟着压。} Word 的 \texttt{w:w="90"} 只作用在字面上，
% 字距是另加的。整行一起缩会把十道 2\,bp 的缝也缩成 1.8，一行差 2\,bp。
%
% \emph{页首那一段的段前要自己留住。} \hologo{TeX} 把换页处的竖胶丢掉，
% 而 Word 在页首照进段前，原件那一行的基线量出 99.12（版心顶 72 加 12 加
% 19\,bp 固定值行盒里基线那一处 15.2）。先发一条零高的 \cs{hrule} 把这一页
% 垫成非空，后面那段胶就留得住。
%
% \cs{topskip} 要一并归零。它管的是页首到首个盒子基线的距离，不归零就再多
% 12.18\,bp——Word 没有这一档，页首那一段与别处一个算法。还回去的时机交给
% \cs{hit@layout@pop@part:}：它里面的 \cs{restoregeometry} 自带一次
% \cs{clearpage}，这一页在那儿排定。
%
% \emph{别自己再发一次 \cs{clearpage}。} 发了会多出一张空白页——底类是双面的，
% 封面那一页排定之后再来一次换页，目录就被推到下一个奇数页去了。
%
% \emph{六段的竖距走类里那套段落模型}（\cs{docxlike\_format\_par:nn}），字号、行距、
% 段前段后都在 \file{src/config/hit-layout.dtx} 的 \option{report-cover-*} 里，
% 这一页只说哪一段排什么。先前是手摆 \cs{vspace}，段与段之间还夹着
% \hologo{TeX} 自己的行距胶，题目那一格的基线因此差了二十多 bp。
%
% 表顶离文种那一行的行盒底 49.75\,bp，原件里表是浮动的，锚在表后那个空段上。
%    \begin{macrocode}
    \cs_new_protected:Npn \hit@szb@art:nnn #1#2#3
      {
        \tl_set:Nx \l__hit_szb_term_tl {#3}
        \int_zero:N \l__hit_szb_idx_int
        \tl_map_inline:Nn \l__hit_szb_term_tl
          {
            \int_compare:nT { \l__hit_szb_idx_int > 0 }
              { \skip_horizontal:n {#2} }
            \int_incr:N \l__hit_szb_idx_int
            \bool_if:nTF {#1} { \scalebox { 0.9 } [ 1 ] { ##1 } } { ##1 }
          }
      }
    \cs_new_protected:Npn \hit@report@shenzhen@bachelor@title@page
      {
        \hit@layout@push@part:n { cover }
        \thispagestyle { empty }
        \dim_set_eq:NN \l__hit_szb_topskip_dim \topskip
        \dim_zero:N \topskip
        \hrule height \c_zero_dim
        \dim_set:Nn \l__hit_szb_indent_dim
          {
            ( \linewidth - \l__hit_szb_label_dim - \l__hit_szb_value_dim ) / 2
          }
        \docxlike_format_par:nn { report-cover-form }
          {
            \hit@if@option@TF { practice } { \hit@empty@box } { \hit@checked@box }
            \hit@term@titlepage@mark@dissertation
            \skip_horizontal:n { 12bp }
            \hit@if@option@TF { practice } { \hit@checked@box } { \hit@empty@box }
            \hit@term@titlepage@mark@practice
          }
        \docxlike_format_par:nn { report-cover-blank } { \mbox { } }
        \docxlike_format_par:nn { report-cover-university }
          {
            \hit@szb@art:nnn { \c_true_bool } { 2bp }
              { \hit@term@cover@university@shenzhen@bachelor@zh }
          }
        \bool_if:NT \g__hit_en_bool
          {
            \docxlike_format_par:nn { report-cover-university-en }
              { \hit@term@cover@university@shenzhen@bachelor@en }
          }
        \bool_if:NT \g__hit_en_bool
          { \docxlike_format_par:nn { report-cover-blank-mid } { \mbox { } } }
        \docxlike_format_par:nn { report-cover-stage }
          {
            \hit@szb@art:nnn { ! \bool_if_p:N \g__hit_en_bool }
              { \bool_if:NTF \g__hit_en_bool { 4bp } { 2bp } }
              {
                \hit@term@cover@form@line@shenzhen@bachelor@zh
                \hit@term@cover@stage@line@shenzhen@bachelor@zh
              }
          }
        \bool_if:NT \g__hit_en_bool
          {
            \docxlike_format_par:nn { report-cover-stage-en }
              {
                \hit@term@cover@form@line@shenzhen@bachelor@en \space
                \hit@term@cover@stage@line@shenzhen@bachelor@en
              }
          }
        \vspace { 49.75bp }
        \hit@szb@title@row:
        \hit@szb@row:nnn { author }      { hit@info@author@zh }      { hit@info@author@en }
        \hit@szb@row:nnn { student@id }  { hit@info@student@id }    { }
        \hit@szb@row:nnn { affiliation } { hit@info@affiliation@zh } { hit@info@affiliation@en }
        \hit@szb@row:nnn { speciality }  { hit@info@major@zh }      { hit@info@major@en }
        \hit@szb@row:nnn { supervisor }  { hit@info@supervisor@zh }  { hit@info@supervisor@en }
        \hit@szb@date@row:
        \par
        \dim_set_eq:NN \topskip \l__hit_szb_topskip_dim
        \hit@layout@pop@part:
      }
%    \end{macrocode}
%
% \emph{题目那一格。} 与六格同一套，只是至少两行：原件那一格是两段，各自带一串
% 下划线，那一格的说明「题目长可分两行写」说的就是这两段。题目自己折到两行也算
% 够了，不再补第二条空线；再长就接着往下长，下面六格跟着让。
%
% 那一行还写着至少 2251 缇即 112.55\,bp（\texttt{w:trHeight} 没写 \texttt{hRule}
% 就是 at-least）。两行各一个步进共 94，凑不满就在末尾补齐，中间那截夹缝正是
% 留给第三行的。
%
% 标签中文版小二、英文版三号，值两版都是黑体小二。
%
% \emph{装不下时整体压缩没有照搬。} iota 那边表长过空表单的底时，会把每一行的
% 步进匀着减一个量，压到字要碰上一行的线为止，再不够就在纸角发红字。这里没做：
% 报告那一档还没有章级版面，压缩要的那几个量（空表单的底、每行能让出多少）都
% 没有落处。做之前表只会往下长，长过版心就翻页。
%    \begin{macrocode}
    \cs_new_protected:Npn \hit@szb@title@row:
      {
        \hit@szb@line:nnnnnnn { report-cover-title } { 18bp } { 2 }
          {
            \hit@szb@pick:nnn { hit@info@title@first@en }
              { hit@info@title@cover@en } { hit@info@title@cover@zh }
          }
          {
            \heiti \xiaoer
            \hit@szb@label:nnnn { title } { 2 } { 3 } { 16bp }
          }
          { \heiti \xiaoer }
          { \l__hit_szb_value_tl }
        \dim_compare:nNnT
          { \l__hit_szb_lines_int \l__hit_szb_pitch_dim } < { 112.55bp }
          {
            \skip_vertical:n
              { 112.55bp - \l__hit_szb_lines_int \l__hit_szb_pitch_dim }
          }
      }
    \hit@hook@new:nn { report-cover-page-shenzhen-bachelor }
      { \hit@report@shenzhen@bachelor@title@page }
  }
%</report-cover-shenzhen-bachelor>
%    \end{macrocode}
%
% \Finale
